\par\noindent\begin{minipage}{\linewidth}
% !TeX root = ../教材题库_全部条目_转写工作稿.tex
% Source: XJ-B1-C01-SU-P028-Y1232
% 第1章 / 小结与复习 小结与复习 / 教材 p.28 / PDF 第35页
% 已对照原页转写题干及公式；待教师复核；未制作答案
% 原题及相邻原解截图（核对用）：../题目原图/XJ-B1-C01-SU-P028-Y1232-01.webp
% 本文件只在转写完成后作为题干使用；OCR 条目尚不是可靠题干。
\begin{bankitem}{第1章 小结与复习 第3题}{第1章 / 小结与复习 小结与复习 / 教材 p.28 / PDF 第35页}{XJ-B1-C01-SU-P028-Y1232}{已对照原页转写题干及公式；待教师复核；未制作答案}
集合是数学的基本概念，数学概念之间的关系，需要用逻辑用语来表达。学习常用逻辑用语可以更深刻地理解数学内容，更严密地进行数学推理以及更准确地表达数学思想。
\end{bankitem}
\end{minipage}\par\medskip
